%%%%%%%%%%%%%%%%%%%%%%%%%%%%%%%%%%%%%%%%%%%%%%%%%%%%%%%%%%%%%%%%%%%%%%
% Lớp 12 | Phương trình mặt phẳng, đường thẳng, mặt cầu | Phương trình mặt phẳng
% Dạng 2H5?1-1: Câu hỏi lý thuyết
% Số câu: 13  (sinh tự động bởi tach_cau.py; từ nay thầy sửa trực tiếp file này)
%%%%%%%%%%%%%%%%%%%%%%%%%%%%%%%%%%%%%%%%%%%%%%%%%%%%%%%%%%%%%%%%%%%%%%

%%%%%%%%%%%%%%%% TRẮC NGHIỆM 4 PHƯƠNG ÁN %%%%%%%%%%%%%%%%

% nguồn: Phan1_Lop_12_Hinh_hoc_Chuong_5
\begin{ex}%[2H5N1-1]%[ma:4c1dfc97]
	Trong không gian với hệ tọa độ $Oxyz$, cho mặt phẳng $(P)\colon x + y - z + 1 = 0$. Điểm nào sau  đây thuộc mặt phẳng $(P)$?
	\choice 
		{\True $A(1;1;3)$} 
		 {$B(1;1;-3)$}
		{$C(3;1;1)$} 
		{$D(-1;-1;3)$} 
	\loigiai{ 
		Ta có $1+1-3+1 = 0  \Rightarrow A(1;1;3) \in (P)$.
	} 
\end{ex}

% nguồn: Phan1_Lop_12_Hinh_hoc_Chuong_5
\begin{ex}%[2H5N1-1]%[ma:6938d9f4]
	Trong không gian $Oxyz$, khoảng cách từ điểm $M\left(x_0;y_0;z_0\right)$ đến mặt phẳng $(P)\colon ax+by+cz+d=0$  bằng
	\choice 
		 {$\dfrac{\left|ax_0+by_0+c\right|}{\sqrt{a^2+b^2}}$}
		{\True $\dfrac{\left|ax_0+by_0+cz_0+d\right|}{\sqrt{a^2+b^2+c^2}}$} 
		 {$\dfrac{\sqrt{a^2+b^2+c^2}}{\left|ax_0+by_0+cz_0+d\right|}$}
		{$\dfrac{\left|ax_0+by_0+cz_0+d\right|}{a^2+b^2+c^2}$} 
	 \loigiai{
		Ta có khoảng cách  từ điểm $M\left(x_0; y_0; z_0\right)$ đến mặt phẳng $(P)\colon ax+by+cz+d=0$ là \[\mathrm{d}\left(M,(P)\right)=\dfrac{\left|ax_0+by_0+cz_0+d\right|}{\sqrt{a^2+b^2+c^2}}.\]
	 }
\end{ex}

% nguồn: Phan1_Lop_12_Hinh_hoc_Chuong_5
\begin{ex}%[2H5N1-1]%[ma:734343b7]
	Trong không gian với hệ tọa độ $Oxyz$,  cho mặt phẳng $(P)\colon x - 2y + z - 5 = 0$. Điểm nào dưới đây thuộc $(P)$?
	\choice 
		 {$P(0;0;-5)$}
		 {\True $M(1;1;6)$}
		{$Q(2;-1;5)$} 
		{$N(-5;0;0)$} 
	\loigiai{ 
		Ta có $1 - 2 \cdot 1 + 6  - 5 = 0$ nên $M(1;1;6)$ thuộc mặt phẳng $(P)$.
	} 
\end{ex}

% nguồn: Phan1_Lop_12_Hinh_hoc_Chuong_5
\begin{ex}%[2H5N1-1]%[ma:07a6402d]
	Trong không gian với hệ toạ độ $Oxyz$, phương trình nào sau đây là một phương  trình tổng quát của mặt phẳng?
	 \choice
		 {$x^2-4y+5z-6=0$}
		{$3x-y^4+5z-6=0$} 
		{\True $3x-4y+5z-6=0$} 
		{$3x-4y+z^5-6=0$} 
	\loigiai{ 
		Trong  các phương trình trên, phương trình có dạng phương trình tổng quát của mặt phẳng là $3x-4y+5z-6=0$.
	} 
\end{ex}

% nguồn: Phan1_Lop_12_Hinh_hoc_Chuong_5
\begin{ex}%[2H5N1-1]%[ma:15cbd2e9]
	Trong không gian với hệ tọa độ $Oxyz$, cho mặt phẳng $(P)\colon 3x - 2z + 2 =  0$ đi qua điểm nào sau đây?
	\choice 
		{$B(4;2;1)$} 
		 {$A(1;2;4)$}
		 {\True $D(2;1;4)$}
		 {$C(2;4;-1)$}
	 \loigiai{
		Thay tọa độ $B(4;2;1)$  vào phương trình $(P)$ ta được
		$3 \cdot 4 - 2 \cdot 1 + 2 =  0 \Leftrightarrow 12 = 0$ (vô lí) $\Rightarrow B \notin (P)$.\\
		Thay tọa  độ $A(1;2;4)$ vào phương trình $(P)$ ta được
		$3 \cdot 1 - 2 \cdot 4 + 2 = 0 \Leftrightarrow  -3 = 0$ (vô lí) $\Rightarrow A \notin (P)$.\\
		Thay tọa độ  $D(2;1;4)$ vào phương trình $(P)$ ta được
		$3 \cdot 2 - 2 \cdot 4 + 2 = 0 \Leftrightarrow 0 = 0$ (thỏa mãn) $\Rightarrow D  \in (P)$.\\
		 Thay tọa độ $C(2;4;-1)$ vào phương trình $(P)$ ta được
		$3  \cdot 2 - 2 \cdot (-1) + 2 = 0 \Leftrightarrow 10 = 0$ (vô lí) $\Rightarrow C \notin (P)$.
	 }
\end{ex}

% nguồn: Phan1_Lop_12_Hinh_hoc_Chuong_5
\begin{ex}%[2H5N1-1]%[ma:9c032141]
	Trong không gian $Oxyz$, cho mặt phẳng $(P)\colon x-2 y+z-5=0$. Điểm nào dưới đây không thuộc  $(P)$?
	 \choice
		 {\True $A(2;-1;-5)$}
		{$B(0;  0; 5)$}
		{$C(5; 0;  0)$}
		{$D(1; 1;  6)$}
	\loigiai{ 
		Thay tọa độ điểm $A(2;-1;-5)$ vào phương  trình mặt phẳng
		$$ 2 - 2(-1) + (-5) - 5 = -6  \neq 0. $$
		Vậy điểm $A$ không thuộc mặt phẳng $(P)$. 
	} 
\end{ex}

% nguồn: Phan1_Lop_12_Hinh_hoc_Chuong_5
\begin{ex}%[2H5N1-1]%[ma:5fa3eefa]
	Trong không gian với hệ tọa độ $Oxyz$, điểm nào sau  đây thuộc mặt phẳng $(Oxy)$?
	 \choice
		{\True  $A(1;1;0)$}
		{$B(1;0;-3)$} 
		 {$C(0;1;1)$}
		{$D(0;0;3)$} 
	 \loigiai{
		Vì mặt phẳng $(Oxy)\colon z =  0$ nên $A(1;1;0)\in (Oxy)$.
	} 
\end{ex}

% nguồn: Phan1_Lop_12_Hinh_hoc_Chuong_5
\begin{ex}%[2H5N1-1]%[ma:abbfbf8b]
	Trong không gian $Oxyz$, phương trình  nào trong các phương trình sau là phương trình tổng quát của một mặt phẳng?
	\choice 
		{$x-y+\dfrac{1}{z}=0$} 
		{\True$x-y-z+2=0$} 
		 {$x+2y+yz-3=0$}
		 {$\dfrac{x}{2}+\dfrac{y}{3}+\dfrac{z}{-1}=1$}
	 \loigiai{
		Phương trình tổng quát của mặt phẳng có dạng $Ax+By+Cz+D=0$  nên $x-y-z+2=0$ là phương trình mặt phẳng.
	 }
\end{ex}

% nguồn: Phan1_Lop_12_Hinh_hoc_Chuong_5
\begin{ex}%[2H5N1-1]%[ma:c85c577a]
	Trong không gian với hệ trục $Oxyz$, cho mặt phẳng $(P)\colon 2x  + y - z + 3 = 0$. Điểm nào sau đây không thuộc mặt phẳng $(P)$?
	 \choice
		{$Q(0; -1;  2)$}
		 {$M(1; -4; 1)$}
		{$P(-1;  0; 1)$}
		 {\True $N(-2; 0; 1)$}
	\loigiai{ 
		Thay từng điểm vào phương  trình mặt phẳng
		\begin{itemize}
			\item Với $Q\colon 2\cdot 0 +  (-1) - 2 + 3 = 0 \Rightarrow$ thuộc mặt phẳng.
			\item Với $M\colon 2\cdot 1 + (-4) - 1 + 3  = 0 \Rightarrow$ thuộc mặt phẳng.
			\item Với $P\colon 2\cdot (-1) + 0 - 1 + 3 = -2 -1 + 3 = 0  \Rightarrow$ thuộc mặt phẳng.
			\item Với $N\colon 2\cdot (-2) +  0 - 1 + 3 = -4 -1 + 3 = -2 \ne 0$ $\Rightarrow$ \textbf{Không thuộc}.
		\end{itemize}
	} 
\end{ex}

% nguồn: Phan1_Lop_12_Hinh_hoc_Chuong_5
\begin{ex}%[2H5N1-1]%[ma:eab9ff1f]
	Trong không gian  $Oxyz$, phương trình nào sau đây \textbf{không} phải là phương trình tổng quát của một mặt phẳng?
	 \choice
		{\True $x - 3y  + z^2 = 0$}
		{$z  - 1 = 0$}
		{$6x - y  + 3z - 1 = 0$}
		{$2x -  z = 0$}
	\loigiai{Phương trình mặt phẳng là phương trình có dạng $Ax+By+Cz+D=0$ trong đó $A$, $B$, $C$, $D$ không  đồng thời bằng $0$.\\
	Phương trình $x - 3y + z^2 = 0$ có $z^2$ là bậc hai nên không là  phương trình của mặt phẳng.
	 }
\end{ex}

% nguồn: Phan1_Lop_12_Hinh_hoc_Chuong_5
\begin{ex}%[2H5N1-1]%[ma:7bcc28c7]
	Trong không gian với hệ trục tọa độ $Oxyz$, cho mặt phẳng  $(\alpha)\colon 2x-3y-z-1=0$. Điểm nào dưới đây \textbf{không} thuộc mặt phẳng $(\alpha)$?
	 \choice
		{$M(-2;1;-8)$} 
		 {$Q(1;2;-5)$}
		{$N(4;2;1)$} 
		 {\True $P(3;1;3)$}
	 \loigiai{
		Thay tọa độ các điểm  vào phương trình mặt phẳng $(P)$.\\
		Ta thấy khi thay tọa độ điểm $P$ thì $2\cdot 3 - 3\cdot 1 - 3 - 1= -1  \ne 0$.\\
		Suy  ra $P\not\in (\alpha)$.
	} 
\end{ex}

% nguồn: Phan1_Lop_12_Hinh_hoc_Chuong_5
\begin{ex}%[2H5H1-1]%[ma:b6da822d]
	Trong không gian $Oxyz$, cho mặt phẳng có phương  trình $2x+y-4z-29=0$. Một điểm thuộc mặt phẳng đã cho~là
	\choice 
		{\True  $(2;1;-6)$}
		 {$(2;1;6)$}
		{$(2;1;-4)$} 
		 {$(-2;-1;8)$}
	\loigiai{ 
		Ta kiểm tra các điểm ở đáp  án.
		\begin{itemize}
			\item $ (2;1;-6) \in 2x+y-4z-29$ vì $2\cdot2+1-4\cdot(-6)-29=0$. 
			\item $ (2;1;6) \notin 2x+y-4z-29$ vì  $2\cdot2+1-4\cdot(-6)-29=-48$.
			\item $ (2;1;-4) \notin  2x+y-4z-29$ vì $2\cdot2+1-4\cdot(-6)-29=-8$.
			\item $ (-2;-1;8) \notin 2x+y-4z-29$ vì $2\cdot2+1-4\cdot(-6)-29=-66$. 
		\end{itemize}
	 }
\end{ex}


%%%%%%%%%%%%%%%% TRẢ LỜI NGẮN %%%%%%%%%%%%%%%%

% nguồn: Phan2_Lop_12_Hinh_hoc_Chuong_5
\begin{bt}%[2H5N1-1]%[ma:57641a97]
	Trong không gian với hệ toạ độ $Oxyz$, cho mặt phẳng $(P)\colon x-2y+2z-5=0$. Điểm $M(5,1,c)$ nằm trên mặt phẳng  $(P)$. Tìm $c$.
	 \shortans{1}
	\loigiai{ 
		Vì điểm $M(5,1,c)$ nằm trên mặt phẳng  $(P)$ nên tọa độ của $M$ thỏa mãn phương trình của $(P)$.\\
		Ta có  $5-2\cdot1+2c-5=0 \Leftrightarrow c=1$.
	 }
\end{bt}
